\section{Related Work}
\label{sec:literature}

\paragraph{Recursive / Looped Transformers.}
\vspace{-1mm}
Parameter sharing provides an orthogonal route to efficiency and effective depth. The Universal Transformer demonstrated that repeatedly applying a single block can match the representational capacity of deep non–shared stacks, while also introducing adaptive computation time \cite{dehghani2018universal}. ALBERT further showed that extensive cross-layer weight-tying yields substantial parameter savings during pretraining, without compromising downstream performance \cite{lan2019albert}. Deep Equilibrium Models (DEQ) extend this paradigm by defining an implicit infinitely deep, weight–tied transformation solved via fixed-point iteration and implicit differentiation \cite{bai2019deep}, with related equilibrium and recurrent architectures also explored in vision and multimodal contexts \cite{jaegle2021perceiver, yang2022recurring}. More recent work has studied looping as programmable computation \cite{giannou2023looped}, iterative data-fitting or optimization solvers \cite{yang2023looped}, mechanisms for algorithmic length generalization \cite{fan2024looped}, and even the Turing-completeness of looped decoders for certain graph algorithms \cite{de2024simulation}. Most relevant to our work, \emph{Time-Modulated Looped Transformers} (TMLT) \cite{xu2024expressive} analyze the expressive power of looped Transformers in language modeling, showing that conditioning on timestep (loop index) improves scaling and perplexity. Our approach builds on this insight but extends it by conditioning on both normalized time and step size, and by training across families of trajectories rather than a single fixed schedule.


\paragraph{Dynamic compute.}
Dynamic compute allocation reduces inference cost by skipping or reallocating computation where it is not needed \cite{bengio2015conditional, huang2016deep, panda2016conditional}. Early exiting halts processing for “easy” inputs at intermediate layers, allowing deeper computation only for harder cases \cite{elbayad2019depth, schuster2022confident, elhoushi2024layerskip}. Other approaches include layer dropping and pruning strategies for BERT-style models \cite{fan2019reducing, hou2020dynabert}. Mixture-of-Experts (MoE) increases capacity through sparse expert routing \cite{shazeer2017outrageously, fedus2022switch}, while \emph{Mixture-of-Depths} (MoD) \cite{raposo2024mixture} reframes adaptivity as token-wise routing across layers, enabling fine-grained dynamic allocation of depth. Despite these advances, the adaptation of dynamic compute techniques to exploit the looping capability of recursive models has received little attention. A concurrent line of work, \emph{Mixture of Recursions} \cite{bae2025mixture}, extends routing to recursive stacks by varying the number of shared-block applications per token. In contrast, our approach does not route or halt computation at the token level; instead, we train looped models to be \emph{budget-conditioned}, ensuring that the same parameters operate robustly across user-specified loop budgets.

\paragraph{Latent reasoning.}
An increasing body of work investigates reasoning carried out within hidden states rather than through explicit chain-of-thought prompting \cite{goyal2023think, cheng2024compressed, pfau2024let, kaissis2026stepresolveddataattributionlooped, chen2026loop,zhu2025scaling, jolicoeur2025less, wang2025hierarchical, zhang2025recursive, knupp2026depth, liang2026latent}. Several approaches employ looped models to simulate multi-step reasoning or to approximate chain-of-thought dynamics directly \cite{hao2024training, saunshi2025reasoning, wang2025hierarchical}. In parallel, theoretical and empirical studies connect network depth and algorithmic generalization to reasoning ability \cite{merrill2023expressive, chen2024can, ye2024physics, xu2025cot}. More recently, analyses suggest that looped Transformers possess an inductive bias for reasoning that strengthens with increasing effective computational depth \cite{saunshi2024inductive, saunshi2025reasoning}. Collectively, this line of work frames such abilities as \emph{latent reasoning}, wherein models perform iterative computations within their hidden-state space without explicit verbalization. Our work leverages this inductive bias of looped models but emphasizes \emph{trajectory robustness}: ensuring that hidden-state computation remains informative even under shorter budgets, while continuing to refine effectively at greater depths.

\paragraph{Time / shortcut modulation and conditioning.}
Conditioning networks on continuous \emph{time} or related control variables has proven highly effective in generative diffusion modeling \cite{ho2020denoising, rombach2022high, lipman2022flow}. Diffusion Transformers (DiT) \cite{peebles2023scalable} incorporate timestep-conditioned adaptive normalization, while consistency-based training aligns solutions across discretizations \cite{song2023consistency}. Recent shortcut and one-step diffusion approaches further distill long trajectories into a few steps by enforcing consistency between coarse and fine solvers \cite{frans2024one, lu2024simplifying}. These ideas are increasingly migrating into language modeling: for example, \emph{Time-Modulated Looped Transformers} (TMLT) \cite{xu2024expressive} adapt DiT-style timestep conditioning to recursive language models, demonstrating improved scaling and perplexity. LoopFormer extends this trend by conditioning each loop not only on normalized time $t$ but also on the step size $\Delta t$, and by training across families of sampled trajectories with a shortcut-consistency objective. This design enables robust performance under arbitrary user-specified compute budgets, without relying on per-token halting or routing mechanisms.

% \paragraph{Transformer ODEs}
% Conceptually our method can be thought of as a combination of DEQ and ODE where we explicitly discretize and condition on the steps as inspired by the diffusion literature. Specially, different from a fixed point solution, like DEQ, we have an online maximum rollout of the model to estimate this point which is hence trainable and flexible, and different from NODE


% \hl{I think the idea is to go through the few ODE and Transformer ODE approaches and say similarities and differences}

